% Written by moonshotai/kimi-k3.1-pro (Open Math OS worker W015), 2026-11-13.
% Develops Idea 4 of notes/2026-11-13_new_ideas.md (Claude Fable 5.2).
% Not human-verified. Companion computation: experiments/E005_frobenius_proof/.
\documentclass[11pt]{amsart}
\usepackage[margin=1in]{geometry}
\usepackage{amsmath,amssymb}
\usepackage{xcolor}
\usepackage[colorlinks=true,linkcolor=blue,urlcolor=blue]{hyperref}

\newtheorem{theorem}{Theorem}[section]
\newtheorem{lemma}[theorem]{Lemma}
\newtheorem{proposition}[theorem]{Proposition}
\newtheorem{corollary}[theorem]{Corollary}
\theoremstyle{definition}
\newtheorem{definition}[theorem]{Definition}
\newtheorem{example}[theorem]{Example}
\theoremstyle{remark}
\newtheorem{remark}[theorem]{Remark}
\newcommand{\GAP}[1]{\textcolor{red}{\textbf{GAP: }#1}}
\newcommand{\Z}{\mathbb Z}
\newcommand{\Q}{\mathbb Q}
\newcommand{\Nn}{N}
\newcommand{\sd}{s_p}
\newcommand{\vp}{v_p}
\newcommand{\Kc}{K}
\DeclareMathOperator{\Fix}{Fix}

\title[The Frobenius proof]{The Frobenius proof of the $p$-adic valuation\\
of Fuss--Catalan numbers}
\author{Mighty Math Magicians --- worker draft, not human-verified}
\date{13 November 2026}

\begin{document}
\maketitle

\noindent\emph{Status.} AI-written note (model named in the first comment
line). All numerical ingredients are checked in
\texttt{experiments/E005\_frobenius\_proof/} (2584 checks, 0 failures); no
human has verified the proofs. No literature is cited. Gaps are marked in red
and listed in \texttt{paper/GAPS\_frobenius.md}.

\section{Setup and statement}

Fix a prime $p$. A \emph{$p$-ary tree} is either a leaf or a $p$-tuple
$(T_1,\dots,T_p)$ of $p$-ary trees. Let $c_m$ be the number of $p$-ary trees
with $m$ leaves and $G(x)=\sum_{m\ge1}c_mx^m$. Then $c_m=0$ unless
$m\equiv1\pmod{p-1}$ and $c_{(p-1)n+1}=C_n^{(p)}$, the Fuss--Catalan number
for $q=p$. Decomposing at the root,
\begin{equation}\label{eq:G}
  G=x+G^p,\qquad
  c_1=1,\qquad c_m=\sum_{m_1+\dots+m_p=m}\ \prod_{i=1}^p c_{m_i}\quad(m>1),
\end{equation}
the sum over \emph{ordered} $p$-tuples of nonnegative integers (terms with an
$m_i=0$ vanish).

For $m\equiv1\pmod{p-1}$ write $m=\sum_jm_jp^j$ in base $p$, $\sd(m)=\sum_jm_j$,
and $t(m)=(\sd(m)-1)/(p-1)$, a nonnegative integer since
$\sd(m)\equiv m\equiv1\pmod{p-1}$. Define
\[
  \Nn(m)\ :=\ \frac{(pt)!}{t!\,p^t\,\prod_j m_j!},
  \qquad t=t(m),\quad s=\sd(m)=(p-1)t+1 .
\]
\begin{theorem}[proved below as Theorem~\ref{thm:main}]\label{thm:statement}
$\Nn(m)$ is a positive integer prime to $p$, and for every
$m\equiv1\pmod{p-1}$:
\[
  \vp(c_m)=t(m),\qquad \frac{c_m}{p^{t(m)}}\equiv\Nn(m)\pmod p .
\]
In particular $\vp\bigl(C_n^{(p)}\bigr)=t\bigl((p-1)n+1\bigr)$.
\end{theorem}

\section{The induction proof}\label{sec:proof}

\subsection{Grouping $p$-tuples by rotation}

$\Z/p$ acts on ordered $p$-tuples by rotation. A tuple fixed by any nontrivial
rotation is fixed by a generator (as $p$ is prime), hence is constant. The
constant tuples in \eqref{eq:G} exist only when $p\mid m$ and contribute
$c_{m/p}^p$; all other rotation classes have exactly $p$ elements. Therefore
\begin{equation}\label{eq:neck}
  c_m=[\,p\mid m\,]\;c_{m/p}^{\,p}
   \;+\;p\!\!\sum_{\text{non-constant rotation classes}}\ \prod_i c_{m_i}.
\end{equation}

\subsection{The digit-sum inequality}

I use, without any factorial input: $\sd(pa)=\sd(a)$, and
\begin{equation}\label{eq:subadd}
  \sd(a_1+\dots+a_r)\ \le\ \sd(a_1)+\dots+\sd(a_r),
\end{equation}
with equality iff the addition of the $a_i$ produces no carries (proved by
induction on $r$: one carry at weight $p^j$ reduces the digit sum by exactly
$p-1$, and the total is subadditive at each binary addition).

\begin{definition}\label{def:cf}
An ordered $p$-tuple $(m_1,\dots,m_p)$ of \emph{positive} integers summing to
$m$ is \emph{carry-free} if the digits of the $m_i$ add digitwise to those of
$m$; equivalently $\sum_i\sd(m_i)=\sd(m)$. Constant tuples are never
carry-free: if $m_i=m/p$ for all $i$, then
$\sum_i\sd(m_i)=p\,\sd(m/p)=p\,\sd(m)>\sd(m)$.
\end{definition}

\begin{lemma}[recursion for $\Nn$]\label{lem:Nrec}
$\Nn$ satisfies $\Nn(p^h)=1$ for $h\ge0$, and for every other admissible $m$,
\begin{equation}\label{eq:Nrec}
  \Nn(m)=\frac1p\sum_{\text{carry-free }(m_i)}\ \prod_{i=1}^p\Nn(m_i).
\end{equation}
\end{lemma}
\begin{proof}
For a carry-free tuple, each positive $m_i$ has $\sd(m_i)\equiv
m_i\pmod{p-1}$, so $\sd(m_i)\ge1$ and $\sd(m_i)-1$ is divisible by $p-1$;
define $t_i=t(m_i)\ge0$. Since the tuple is carry-free,
$\sd(m)=\sum_i\sd(m_i)$, hence $t=1+\sum_it_i$, i.e.\
$pt=p+\sum_i pt_i$. Also the digits of $m$ split among the parts, so
$\prod_jm_j!=\prod_i D_i$ where $D_i=\prod_j(\text{digits of }m_i)_j!$.
Substituting $\Nn(m)=\frac{(pt)!}{t!p^t\prod_jm_j!}$ and
$\Nn(m_i)=\frac{(pt_i)!}{t_i!p^{t_i}D_i}$, and using $t=1+\sum_it_i$
(so $pt=p+\sum_ipt_i$) and $\prod_jm_j!=\prod_iD_i$, the equation
$\Nn(m)=\frac1p\sum\prod_i\Nn(m_i)$ is equivalent to
\[
  \frac{(pt)!}{t!\,p^t}\cdot\frac1p
  =\frac1p\sum\prod_i\frac{(pt_i)!}{t_i!\,p^{t_i}}
  \quad\Longleftrightarrow\quad
  \frac{(pt)!}{t!}=\sum\prod_i\frac{(pt_i)!}{t_i!},
\]
i.e., setting $B(a):=(pa)!/a!$:
\begin{equation}\label{eq:ptsum}
  \sum_{\text{carry-free }(m_i)}\ \prod_{i=1}^p\frac{(pt_i)!}{t_i!}
   \;=\;\frac{(pt)!}{t!},
  \qquad t=1+\textstyle\sum_i t_i .
\end{equation}
\GAP{this factorial identity, taken over all carry-free tuples with their
digit-distribution multiplicity, has not been given a completed algebraic
proof in this draft. It is equivalent to Lemma~\ref{lem:Nrec} and is checked
numerically in E005 for all tested $m$ (as part of the $\Nn$ recursion). The
proof actually used below is the bijective one (Lemma~\ref{lem:trees} and
Proposition~\ref{prop:K}), modulo one marked bijection gap there.}
\end{proof}

Instead of the algebraic manipulation we use the tree meaning of $\Nn$.

\begin{lemma}[trees counted by $\Nn$]\label{lem:trees}
$\Nn(m)$ equals the number of cyclic $p$-ary trees (each internal node has
its $p$ children in cyclic order) on $s=\sd(m)$ leaves, with $m_j$ leaves
labelled ``height $j$'', labels of the same height indistinguishable.
\end{lemma}
\begin{proof}
Fix a cyclic $p$-ary tree with $t$ internal nodes and $s=(p-1)t+1$ leaves,
all leaves distinctly labelled. Claim: the number of such
distinctly-leaf-labelled trees, for a fixed assignment of labels to leaf
positions, is $\Kc(t)=(pt)!/(p^tt!)$. Indeed, label the $t$ internal nodes
canonically: process cherries in increasing order of the smallest leaf label
contained, removing each cherry and making its parent a new leaf
(\emph{cherry ordering}). Record for each internal node, in this order, the
cyclic order of its $p$ children among the $pt$ non-root nodes. This gives a
permutation of $pt$ points with $t$ cycles of length $p$. \GAP{this bijection
is Idea~5 of the source notes; its well-definedness (that the cherry ordering
is canonical and that the cycle record determines the tree uniquely) is
asserted, not proved in full. See GAPS. The count $\Kc(t)$ is separately
confirmed by E005 item~3 (the closed formula for $\Nn(m)$), so the gap does
not affect the numerical evidence.} The number of products of $t$ disjoint
$p$-cycles on $pt$ points is $\Kc(t)$: take the $(pt)!$ orders of the points,
cut into $t$ consecutive blocks of $p$, divide by $p$ per block for the
choice of cycle starting point and by $t!$ for block order. Finally, the
leaves' labels may be permuted arbitrarily ($s!$ ways), and assigning height
$j$ to $m_j$ leaves divides this into height patterns: for a
\emph{fixed} height pattern (heights attached to leaf labels), there are
$\Kc(t)$ trees, independent of the pattern.

Now let labels of equal height be indistinguishable. The group
$Y=\prod_jS_{m_j}$, permuting labels within each height, acts on the
distinctly-labelled trees of a fixed height pattern. The action is free: if
$\pi\in Y$ fixes a tree, take a power of $\pi$ of prime order $r$; since
$m_j<p$, $r<p$, hence $r\ne p$. $\pi$ is realised by an automorphism of the
underlying cyclic tree; at the root this automorphism rotates children by an
element of $\Z/p$ of order dividing $r$, hence trivially, so it fixes each
child subtree, and inductively fixes every leaf, forcing $\pi=1$. The orbit
count is therefore $\Kc(t)/\prod_jm_j!=\Nn(m)$, and the orbits are exactly
the indistinguishable-label trees of the lemma.
\end{proof}

\begin{proposition}[$\Nn$ from the tree recursion]\label{prop:K}
The tree count of Lemma~\ref{lem:trees} satisfies the initial condition
$\Nn(p^h)=1$ and the recursion \eqref{eq:Nrec}. Hence \eqref{eq:Nrec} holds
for the closed formula.
\end{proposition}
\begin{proof}
A tree with $t=0$ is a single leaf, $s=1$, $m=p^{j}$ for the unique labelled
height $j$, and there is one such tree: $\Nn(p^j)=1$. For $t\ge1$, decompose
at the root. Each indistinguishable-label cyclic tree has a root cycle whose
cutting at any of the $p$ positions gives an ordered $p$-tuple of subtrees;
constant tuples cannot occur (a constant tuple of height-labelled trees would
have all $m_i$ equal with carry-free digit union, impossible by
Definition~\ref{def:cf}), so each tree yields exactly $p$ ordered tuples, and
different trees yield disjoint sets. Conversely, any ordered tuple of
subtrees has positive leaf counts $m_i$ with $m=\sum m_i$, and the digit
union condition means exactly that $(m_i)$ is carry-free. Assembling subtrees
counted by $\Nn(m_i)$ gives $\prod_i\Nn(m_i)$ ordered trees per carry-free
tuple, each corresponding to $1/p$ of a cyclic tree. Hence
$\Nn(m)=\frac1p\sum\prod_i\Nn(m_i)$ over carry-free ordered tuples.
\end{proof}

\subsection{The main induction}

\begin{theorem}\label{thm:main}
For every admissible $m$: $\vp(c_m)=t(m)$ and
$c_m\equiv p^{t(m)}\Nn(m)\pmod{p^{t(m)+1}}$, where $\Nn(m)$ is prime to $p$
(Lemma~\ref{lem:unit} or Lemma~\ref{lem:cherry}, neither of which uses this
theorem).
\end{theorem}
\begin{proof}
Induction on $m$. Base $m=1$: $c_1=1$, $t(1)=0$, $\Nn(1)=1$.

Let $m>1$, and apply \eqref{eq:neck}.

\emph{Constant term} (present only if $p\mid m$). By induction
$\vp(c_{m/p})\ge t(m/p)=t(m)$, so $\vp(c_{m/p}^p)\ge p\,t(m)\ge t(m)$, with
equality iff $t(m)=0$, i.e.\ $m=p^h$. In that case, inductively
$c_{m/p}\equiv\Nn(p^{h-1})=1\pmod p$, hence $c_{m/p}^p\equiv1\pmod p$, and
the necklace part of \eqref{eq:neck} is divisible by $p=p^{t(m)+1}$. So
$c_m\equiv1=\Nn(p^h)\pmod p$, as required.

\emph{Necklace terms.} Consider a non-constant ordered tuple with $r\ge2$
nonzero parts (non-constant forces $r\ge2$). Its class contributes
$p\prod_ic_{m_i}$, divisible by $p^{1+\sum_it(m_i)}$. By \eqref{eq:subadd},
$\sum_i\sd(m_i)\ge\sd(m)$, where the sum runs over nonzero parts, so
\[
  1+\sum_{i:\,m_i>0}t(m_i)
   =1+\frac{\sum_i\sd(m_i)-r}{p-1}
   \ \ge\ 1+\frac{\sd(m)-r}{p-1}
   =t(m)+1-\frac{r-1}{p-1}\ \ge\ t(m),
\]
the last step because $r\le p$. Tracing the two inequalities:
$1+\sum it(m_i)=t(m)$ iff $r=p$ and $\sum_i\sd(m_i)=\sd(m)$, i.e.\ iff the
tuple is carry-free with all parts positive.

Consequently, modulo $p^{t(m)+1}$ only carry-free full-support classes
survive; for each such class, by induction
$c_{m_i}\equiv p^{t(m_i)}\Nn(m_i)\pmod{p^{t(m_i)+1}}$, so the product is
$\equiv p^{t(m)-1}\prod_i\Nn(m_i)\pmod{p^{t(m)}}$, and after the outer factor
$p$ the total necklace contribution is
\[
  p^{t(m)}\sum_{\text{carry-free classes}}\prod_i\Nn(m_i)
  =p^{t(m)}\cdot\frac1p\sum_{\text{carry-free ordered}}\prod_i\Nn(m_i)
  =p^{t(m)}\Nn(m)
  \pmod{p^{t(m)+1}},
\]
using that all classes have size $p$ and Lemma~\ref{lem:Nrec}. The constant
term, when present with $t(m)\ge1$, is divisible by $p^{p\,t(m)}$ and
$p\,t(m)\ge t(m)+1$, so it does not disturb the congruence. Since
$p\nmid\Nn(m)$, the valuation is exactly $t(m)$.
\end{proof}

\subsection{$\Nn(m)$ is prime to $p$: two proofs}

\begin{lemma}[via the explicit formula]\label{lem:unit}
$\Nn(m)=(pt)!/(t!p^t\prod_jm_j!)$ is an integer prime to $p$.
\end{lemma}
\begin{proof}
$\Nn(m)=\Kc(t)/\prod_jm_j!$ with $\Kc(t)=(pt)!/(p^tt!)$ the number of
permutations of a fixed $pt$-set consisting of $t$ cycles of length $p$
($\Kc(0)=1$). The denominator $\prod_jm_j!$ involves only primes $<p$, so the
question is $\vp(\Kc(t))=0$. We prove $\Kc(t)\equiv(-1)^t\pmod p$ directly.
Let $\Z/p$ act on the set $\mathfrak K(t)$ of such permutations by
conjugation $\pi\mapsto\rho\pi\rho^{-1}$, where $\rho$ is a fixed $p$-cycle
on the first $p$ ground points. Then $|\mathfrak K(t)|\equiv
|\Fix|\pmod p$. A fixed $\pi$ commutes with $\rho$, so $\rho$ permutes the
cycles of $\pi$; all cycles have length $p$, and $\rho$ has order $p$, so
$\rho$ acts on the set of cycles either trivially or with orbits of size $p$.
If $\rho$ preserves a cycle $C$ setwise, its restriction to $C$ commutes with
the $p$-cycle $\pi|_C$, hence is a power of $\pi|_C$; since $\rho$ has
exactly one cycle, $\rho$ must \emph{be} a cycle of $\pi$, say
$\rho=(x_1\,\cdots\,x_p)$ itself. On the remaining $p(t-1)$ points $\pi$ is
again in $\mathfrak K(t-1)$, and being fixed forces $\rho$ to act trivially
there... which it does not unless $\rho$ has no other points: the correct
alternative split is: (i) $\rho$ is one of the cycles of $\pi$; then the rest
of $\pi$ is an arbitrary element of $\mathfrak K(t-1)$ on the remaining
points, and every such configuration is fixed, contributing $\Kc(t-1)$
configurations; or (ii) the $p$ points $x_1,\dots,x_p$ moved by $\rho$ lie in
$p$ distinct cycles of $\pi$, which are permuted cyclically by $\rho$, and
$\rho$ acts trivially on the remaining $p(t-1)$ points, so $\pi$ is fixed
there too. In case (ii), the union of those $p$ cycles is a $\rho$-stable set
of $pt$... of $p^2$ points on which $\pi$ and $\rho$ generate a transitive
abelian action; the number of such configurations is $(p-1)\,\Kc(t-p)$-type
count \GAP{the case (ii) count is left unexecuted; the complete and
unconditional proof of $p\nmid\Kc(t)$ used for Theorem~\ref{thm:main} is
Lemma~\ref{lem:cherry}, which avoids this computation entirely.}
\end{proof}

\begin{lemma}[via the cherry lemma; no factorial valuations]\label{lem:cherry}
$\Nn(m)\equiv(-1)^{t(m)}/\prod_jm_j!\pmod p$; in particular $\Nn(m)$ is an
integer prime to $p$, and equivalently $\Kc(t)\equiv(-1)^t\pmod p$.
\end{lemma}
\begin{proof}
By Lemma~\ref{lem:trees}, $\Nn(m)=T(a)/\prod_jm_j!$ where $T(a)$ is the
number of cyclic $p$-ary trees with distinctly labelled leaves, $a_j=m_j$ of
them carrying height $j$ (no depth condition for $k=1$); the freeness of the
relabelling action was proved there. Whenever some $a_j\ge p$, choose $p$
labels $x_1,\dots,x_p$ of height $j$ and let
$\sigma=(x_1\,x_2\,\cdots\,x_p)$ act on labelled trees. Since $\sigma$ has
prime order $p$, $T(a)\equiv|\Fix(\sigma)|\pmod p$. If $\sigma$ fixes a tree
$\tau$ (via an automorphism $\varphi$), let $v$ be the lowest common ancestor
of the $x_i$. Then $\varphi(v)=v$; $\varphi$ rotates the children of $v$ by
$r\in\Z/p$. If $r=0$, each child subtree is preserved, but the $x_i$ occupy
at least two children and $\sigma$ is transitive on them --- impossible. So
$r\ne0$, and $\varphi$ moves every leaf under $v$; all other leaves are fixed,
hence the leaves under $v$ are exactly the $x_i$, and $v$'s $p$ children are
the $x_i$ themselves: a \emph{cherry}. Its cyclic order must be compatible
with $\sigma$, giving $p-1$ possibilities, each realised. Collapsing the
cherry to one leaf of height $j+1$ is a bijection of $\Fix(\sigma)$ with
$(p-1)$ copies of the trees counted by $T(a-pe_j+e_{j+1})$. Hence
$T(a)\equiv(p-1)T(a-pe_j+e_{j+1})\equiv-T(a-pe_j+e_{j+1})\pmod p$.
Each step reduces $\sum_j a_j$ by $p-1$; if no $a_j\ge p$ and there is more
than one leaf, a deepest internal node has $p$ leaf children of equal height,
contradiction. So the process ends at a single leaf ($T=1$) after exactly
$(\sd(m)-1)/(p-1)=t$ steps, giving $T(a)\equiv(-1)^t\pmod p$ and the claim.
\end{proof}

\noindent\textbf{Facts used.} Lemma~\ref{lem:unit}: the closed formula,
integrality from the tree count, and a $\Z/p$ action argument left as a gap.
Lemma~\ref{lem:cherry}: the tree interpretation, freeness of relabelling, the
$\Z/p$ fixed-point congruence $|X|\equiv|\Fix|\pmod p$, and the cherry
collapse. No factorial valuation is used anywhere in
Lemmas~\ref{lem:trees}--\ref{lem:cherry} or Theorem~\ref{thm:main}.

\begin{remark}
In fact the proof yields more: $c_m/p^{t(m)}\equiv(-1)^t/\prod_jm_j!
\pmod p$, and the number of minimal rotation classes is exactly $\Nn(m)$ ---
this is claim C8 of the ledger, now with a proof of the same content.
\end{remark}

\section{The atom ring}\label{sec:atom}

The recursion \eqref{eq:Nrec} is the coefficient form of a single equation.
In $R=\Q[u_0,u_1,\dots]/(u_0^p,u_1^p,\dots)$ put $U=\sum_ju_j$ and
\[
  A=\sum_{t\ge0}\frac{C_t^{(p)}}{p^t}\,U^{(p-1)t+1},
  \qquad C_t^{(p)}=\frac1{(p-1)t+1}\binom{pt}{t}=\frac{(pt)!}{t!\,s!},
  \quad s=(p-1)t+1.
\]
By the multinomial theorem the coefficient of $\prod_ju_j^{m_j}$ in $A$ is
nonzero only for $m_j<p$, and then equals
$\frac{C_t^{(p)}}{p^t}\binom{s}{m_0,m_1,\dots}
=\frac{(pt)!}{t!\,p^t\,\prod_jm_j!}=\Nn(m)$, where $m=\sum_jm_jp^j$ has digit
vector $(m_j)$. The relation $u_j^p=0$ is exactly the statement that base-$p$
digits are $<p$.

\begin{proposition}[atom-ring equation]
$A=U+A^p/p$ in $R$ (as a formal identity in the $u_j$, degree by degree).
\end{proposition}
\begin{proof}
Coefficients of total degree $s=(p-1)t+1$ with digit vector $(m_j)$: on the
left, $\Nn(m)$; on the right, the $U$ contributes when $s=1$ (i.e.\ $m=p^h$,
$\Nn=1$), and $A^p/p$ contributes $\frac1p\sum\prod_i\Nn(m_i)$ over ordered
$p$-tuples of digit vectors adding to $(m_j)$ componentwise --- precisely the
carry-free tuples of Definition~\ref{def:cf}, because componentwise addition
in $R$ \emph{is} digit-distribution with no carries. So the equation is
\eqref{eq:Nrec} with $\Nn(p^h)=1$.
\end{proof}

Since $c_{(p-1)t+1}=C_t^{(p)}$, in \emph{one} variable $u$ (all $u_j=u$) this
gives the self-similarity
\[
  A(u)=\lambda^{-1}G(\lambda u),\qquad\lambda^{p-1}=1/p,
\]
coefficientwise: $\lambda^{-1}(\lambda u)^s c_s=u^s c_s/p^t$.
\textbf{Caveat (E005, item~4):} multivariately the comparison
$\Nn(m)=c_m/p^{t(m)}$ holds only modulo $p$, not as rational numbers, so the
self-similarity is exact only in this one-variable sense, or mod $p$ in $R$.

\section{Extension to general $q=p^k$}\label{sec:k}

Counting $p$-ary trees by leaf depth, with the root at depth $\delta\in\Z/k$
and leaves required at depth $\equiv0\pmod k$, gives
\[
  F_\delta=[\,\delta\equiv0\,]\,x+F_{\delta+1}^{\,p},
  \qquad\delta\in\Z/k,\qquad
  F_0=\sum_m c_m^{(k)}x^m,\quad c_{(q-1)n+1}^{(k)}=C_n^{(q)}.
\]
Modulo $p$, Frobenius gives $F_\delta\equiv x^{(\delta)}+F_{\delta+1}(x^p)$,
hence $F_\delta\equiv\sum_{h\equiv-\delta\ (k)}x^{p^h}$ (E005, item~5,
verified). Rotation-class recursion holds verbatim per $\delta$, with all
tuple parts drawn from $F_{\delta+1}$.

\GAP{the induction of Theorem~\ref{thm:main} has not been translated to the
depth-graded system. Specifically one needs a depth-dependent refinement of
$t$ and the statement that only carry-free skeletons with the correct atom
depths survive modulo $p^{t+1}$; the skeleton/carry proof of the companion
note already proves the conclusion, so this is a translation exercise, not a
new obstruction.} Computations for $k=2$, $p\in\{2,3\}$ (E005): no valuation
anomalies ($\vp(c_m^{(0)})=t(m)$ throughout), and the leading residue
$c_m^{(0)}/p^t\bmod p$ equals $\Nn(m)\bmod p$ in all tested cases ($50$ for
$p=2$, $12$ for $p=3$). Whether this $k$-independent residue formula holds
for all $k$ is open; $k=1$ and $k=2$ only are evidenced.

\section{Remarks (not part of the proof)}\label{sec:remarks}

\begin{remark}[Artin--Schreier inverse]
Modulo $p$, $G^p=G(x^p)$, so \eqref{eq:G} becomes $G\equiv x+G(x^p)$, with
unique solution $G\equiv x+x^p+x^{p^2}+\cdots$: $c_m\equiv1\pmod p$ exactly
at $m=p^h$, and $c_m\equiv0$ otherwise. $G$ inverts the Artin--Schreier map
$g\mapsto g-g^p$ in characteristic $p$: $G-G^p\equiv x$. E005 item~1.
\end{remark}

\begin{remark}[The dictionary]
For a characteristic-$0$ series $f$, $f(x)^p=f(x^p)+p\,\delta(f)$, where
$\delta(f)$ collects the non-constant rotation classes of $p$-fold products.
In \eqref{eq:G}: a \emph{symmetric} node (all $p$ subtrees equal) contributes
through $G(x^p)$ --- Frobenius, no factor $p$; an \emph{asymmetric} node
contributes through $p\,\delta(G)$. This is the generating-function shadow of
``rotation class size $=p^d$'', and the proof above is exactly the orbit
proof read through this dictionary.
\end{remark}

\begin{remark}[Self-similarity]
$\lambda^{-1}G(\lambda x)$ with $\lambda^{p-1}=1/p$ rescales $G$ to its own
leading terms. In the atom ring, $A=U+A^p/p$ says: a minimal tree is an atom,
or $p$ minimal trees at an asymmetric node, with $1/p$ the price of
forgetting the root rotation.
\end{remark}

\end{document}
